% 本文件由 LaTeX 排版 Agent 根据有效 translation.json 自动生成。
% author=Tatian; work_id=0203; title=塔提安残篇

\chapter{\WorkTitle{塔提安残篇}}

% structure: p0001 source=h1 target=omitted reason=page-title-already-rendered-by-parent

% structure: p0002 source=h2 target=section reason=relative-heading-rank; base=h2

\section{一}

在他的论著\WorkTitle{论救主所说的成全}中，他写道：\Quote{同意确实适合祈祷，但参与败坏则削弱恳求。无论如何，他通过许可确实禁止了，尽管是委婉地；因为他虽然因撒殚和不节制而允许他们回到同一件事上，但他显明了一个将试图事奉两个主人的——藉着\InnerQuote{同意}事奉天主（\BibleRef{格前 7:5}），但藉着不同意，则是不节制、淫乱和魔鬼。}——\Signature{Clem. Alex.}: \WorkTitle{Strom.}, iii. c. 12.\par

% structure: p0004 source=h2 target=section reason=relative-heading-rank; base=h2

\section{二}

某人攻击生育，称之为腐败和毁灭；另有人强解[圣经]，将救主的话应用于生育：\Quote{不要在地上积攒财宝，那里有虫蛀和锈蚀；}并且他不羞于加上先知的话：\Quote{你们都要像衣服一样变旧，蛀虫将吞吃你们。}\par

同样，他们引用关于死人复活的话：\Quote{那世界的儿子们既不娶也不嫁。}——\Signature{Clem. Alex.}: iii. c. 12, § 86.\par

% structure: p0007 source=h2 target=section reason=relative-heading-rank; base=h2

\section{三}

塔提安主张基督虚幻的身体，断定一切性交皆为不洁，他也是厄尼克拉提派极其暴力的异端首领，他使用了这样的论证：\Quote{若有人撒种于肉身，必从肉身收获败坏；}但撒种于肉身的是与女人结合的人；因此，娶妻并在肉身撒种的人，必从肉身收获败坏。——\Signature{Hieron.}: \WorkTitle{Com. in Ep. ad Gal.}\par

% structure: p0009 source=h2 target=section reason=relative-heading-rank; base=h2

\section{四}

他脱离教会，因对其老师的自负而趾高气扬，自以为超越他人，便形成了自己独特的教义类型。他想象某些无形的永世，类似\Person{瓦伦提诺}的，并与\Person{马西昂}和\Person{萨图尼努斯}一样，谴责婚姻为污秽和淫乱，并否认亚当的救恩，这是他自己的意见。——\Signature{Irenæus}: \WorkTitle{Adv. Hœr.}, i. 28.\par

% structure: p0011 source=h2 target=section reason=relative-heading-rank; base=h2

\section{五}

塔提安时而试图引用保禄的话，说在亚当内众人都死了，却忽略了\Quote{罪在那里显多，恩宠却更加显多}。——\Signature{Irenæus}: \WorkTitle{Adv. Heres.}, iii. 37.\par

% structure: p0013 source=h2 target=section reason=relative-heading-rank; base=h2

\section{六}

反对塔提安，他说\Quote{要有光}这些话应被视为祈祷。如果说话的那位知道一位更高的天主，那么祂怎么说\Quote{我是天主，在我以外没有别的天主}呢？\par

他说，对于亵渎、愚妄言谈和放荡的话语，都有惩罚，由圣言处罚和惩戒。他又说，妇女因头发和装饰而受罚，由一位管理这些事的能力所惩罚，这能力也曾藉\Person{参孙}的头发赐予他力量，并惩罚那些因头发的装饰而被引向淫乱的人。——\Signature{Clem. Alex.}: \WorkTitle{Frag.}\par

% structure: p0016 source=h2 target=section reason=relative-heading-rank; base=h2

\section{七}

但塔提安不明白\Quote{要有}这种表达并非总是祈愿式的，有时是命令式的，因此极其不敬地想象那位说\Quote{要有光}的天主是在祈祷，而不是命令有光，好像，正如他不敬地认为，天主是在黑暗中。——\Signature{Origen}: \WorkTitle{De Orat.}\par

% structure: p0018 source=h2 target=section reason=relative-heading-rank; base=h2

\section{八}

塔提安分开旧人和新人，但并非如我们所说，理解旧人是法律，新人是福音。我们同意他说的同一件事，但并非以他所愿的意义，他废除法律，仿佛它属于另一位天主。——\Signature{Clem. Alex.}: \WorkTitle{Strom.}, iii. 12.\par

% structure: p0020 source=h2 target=section reason=relative-heading-rank; base=h2

\section{九}

塔提安不仅谴责和拒绝婚姻，也拒绝天主为使用而创造的食物。——\Signature{Hieron.}: \WorkTitle{Adv. Jovin.}, i. 3.\par

% structure: p0022 source=h2 target=section reason=relative-heading-rank; base=h2

\section{十}

\Quote{但你们将酒给纳齐尔人喝，并命令先知说，不要预言。}关于这一点，或许厄尼克拉提派的首领塔提安试图建立他的异端，主张不可饮酒，因为法律中命令纳齐尔人不可饮酒，而今那些给纳齐尔人酒喝的人被先知指责。——\Signature{Hieron.}: \WorkTitle{Com. in Amos.}\par

% structure: p0024 source=h2 target=section reason=relative-heading-rank; base=h2

\section{十一}

厄尼克拉提派的始祖塔提安，他自己拒绝了一些保禄的书信，尤其认为这封，即[写给]弟铎的，应被宣认为宗徒的，他轻视\Person{马西昂}和其他一些同意他这一点的人的主张。——\Signature{Hieron.}: \WorkTitle{Præf. in Com. ad Tit.}\par

% structure: p0026 source=h2 target=section reason=relative-heading-rank; base=h2

\section{十二}

[阿尔赫劳（主后280年），美索不达米亚卡哈城的主教，将他的同乡塔提安与\Quote{马西昂、撒伯里乌及其他为自己编造了一门特殊科学的人}，即他们自己的神学，列为一类。——\Signature{Routh}: \WorkTitle{Reliquiæ}, tom. v. p. 137. 但参阅本著作的爱丁堡系列，卷二十，第267页。]\par

\SourceRecord{Translated by J.E. Ryland.；Ante-Nicene Fathers；Vol. 2.；Edited by Alexander Roberts, James Donaldson, and A. Cleveland Coxe.；Buffalo, NY: Christian Literature Publishing Co.,；1885.；<http://www.newadvent.org/fathers/0203.htm>.}
